\section{Introduction} 

Developers are increasingly tasked with modifying and extending operating system kernels to improve performance, security, and reliability, or introducing new features to their systems. Extended Berkeley Packet Filter (eBPF) have emerged as the de facto method for extending an operating system, with recent support for both Linux and Windows~\cite{eBPFWindows}.  
eBPF programs inject new logic that is executed before or after existing kernel logic to observe or modify the kernel's behavior. 
eBPF programs were originally used to trace network traffic, but the ecosystem now provides sufficient power to implement a variety of features, including performance monitoring performance~\cite{gregg16,gregg2021computing}, intrusion detection~\cite{Bachl21,jia2023practical}, and application-specific logic~\cite{Zhong22,yang2023lambda,jia2023programmable,ghigoff2021bmc}. eBPF has emerged as a powerful observability framework, traditionally employed in network analytics~\cite{ghigoff2021bmc,chen2025etran,xu2025state,benson2024netedit,zhou2024dint}, system verification~\cite {wu2025vep,bpftime}, system performance\cite{mao2024merlin}, system security~\cite {jin2024beebox,wang2024seak}, and system profiling~\cite{zhong2022xrp,yang2023lambda}. eBPF allows lightweight, sandboxed instrumentation programs to execute in kernel space without intrusive kernel changes or performance penalties. \xiangyu{I think we should simplify the description of eBPF within 1-2 sentences, especially in the intro section.}
However, existing eBPF techniques remain CPU-centric, constrained by frequent kernel-user space context switches and limited visibility into the GPU execution model.

\xiangyu{Currently, }the rapid growth of data-intensive and computationally demanding workloads—ranging from large-scale machine learning and high-performance computing (HPC) simulations~\cite{kousha2019designing,shen2018cudaadvisor,zhou2021measurement,stocks2024breaking,das2023large,fedeli2022pushing,liu2021closing} to graphics rendering~\cite{pankratz2021vulkan,sellers2016vulkan} and large language models~\cite{guan2024amanda,kwon2023efficient}—has increasingly relied on GPU acceleration to achieve desired performance and throughput. 
As these GPU-powered systems scale, precise observability and low-overhead instrumentation become essential for diagnosing performance bottlenecks, optimizing resource usage\cite{hong2017gpu,kato2011timegraph}, and identifying anomalous behavior at runtime. 
Existing GPU instrumentation approaches, such as NVIDIA’s CUDA Profiling Tools Interface (CUPTI)~\cite{cupti} and binary-instrumentation frameworks like NVBit~\cite{villa2019nvbit}, offer insights into GPU performance but often incur substantial overhead or require invasive modifications that interrupt GPU kernels.\xiangyu{Could you cite related work to mention the substantial overhead?}
\xiangyu{eBPF has the potential to monitor GPU performance with low overhead.}

\xiangyu{However, } applying eBPF to GPU workloads presents unique challenges: GPU execution models feature asynchronous kernel launches, parallel thread execution, and specialized memory hierarchies, making instrumentation significantly more complex and prone to high overhead. Furthermore, eBPF’s inherent CPU-centric design does not directly extend to GPU hardware without novel translation and runtime adaptation strategies.

To bridge this gap, we introduce \sys, the first observability framework and eBPF runtime that dynamically extends eBPF instrumentation into GPU kernels via runtime PTX injection. 
Specifically, \sys{} compiles eBPF bytecode into NVIDIA's Parallel Thread Execution (PTX) intermediate representation at runtime and injects this instrumentation directly into actively executing GPU kernels. 
Leveraging user-space eBPF uprobes, shared-memory maps, and real-time GPU telemetry, \sys{} provides fine-grained GPU instrumentation with minimal runtime overhead. 
Unlike previous GPU profiling solutions, our approach \xiangyu{shows more flexibility by} allowing instrumentation to be dynamically added, modified, or removed when running GPU kernels without interrupting their execution. 

\xiangyu{I think we should mention 3 use cases (MemTrace, Fair Share Scheduling, CPU-GPU Collaborative Inference) in the intro section to highlight \sys{}'s versatility.}

% Moreover, while initially demonstrated through precise GPU observability, the broader significance of dynamically injecting eBPF onto GPUs extends beyond instrumentation alone. 
\xiangyu{More broadly,} our methodology opens a window for programmable GPU computing, potentially enabling adaptive runtime kernel optimizations, dynamic GPU resource management, real-time security policy enforcement, and other novel runtime GPU modifications previously inaccessible to existing static or high-overhead instrumentation approaches.% The code is open-sourced under the \url{https://github.com/eunomia-bpf/eGPU} repo.

In summary, we make the following core contributions:
\begin{itemize}[leftmargin=*,itemsep=0pt]
    \item We introduce \sys{}, the first system to dynamically offload eBPF instrumentation and bytecode directly onto running GPU kernels using real-time PTX injection, significantly reducing instrumentation overhead compared to existing methods.
    
    \item We detail the design and implementation of \sys{}, which integrates kernel-level and user-space eBPF instrumentation hooks, runtime PTX generation, and shared-memory synchronization, providing a seamless, low-overhead observability platform for modern HPC and AI workloads. \xiangyu{the first 2 contributions are similar? Every paper would introduce the concept and then implement it. Maybe we should highlight the unique part of \sys{}.}
    
    \item We evaluate \sys{} through initial micro-benchmark experiments measuring GPU memory access latency instrumentation, demonstrating practical low overhead and high-resolution insights into GPU execution. \xiangyu{Do we have number here?}
    
    \item We discuss broader implications and future possibilities enabled by \sys{}, highlighting its potential as a foundational infrastructure for three versatile representative use cases, MemTrace, fair share scheduling, and ML Green. \xiangyu{Should we start the sentence by ``we discuss"? Maybe we should say that \sys{} opens a venue for programmable GPU computing?}

% \begin{itemize}[leftmargin=*,itemsep=0pt]
%   \item \textbf{MemTrace:} Per‐warp memory‐access tracing to detect cache‐miss patterns and bank conflicts inside active kernels, enabling pinpoint diagnosis of memory‐bound bottlenecks.
%   \item \textbf{Fair share Scheduling:} Real‐time collection of sharing information during execution to rebalance workloads and reduce tail latencies in irregular or unbalanced kernels.
%   \item \textbf{CPU–GPU Collaborative Inference (ML Green):} Sub‐100 µs end‐to‐end latency breakdowns across host copy routines and GPU kernels for large‐model pipelines (e.g., Llama 2), driving dynamic batch‐size adaptation and prefetch ordering.
% \end{itemize}
\end{itemize}


